The distinction between a tile and a particular tiling is essential: a tile
may admit both periodic and nonperiodic translation sets.  The theorem
excludes every periodic translation set for the constructed tile.
For lattice tilings, dimension three is the smallest possible dimension
for this conclusion. Newman's finite-state argument shows that the
translation set of every tiling of $\Z$ by a finite tile is periodic
\citep[proof of Theorem~2]{Newman}. Bhattacharya proved that every finite tile of $\Z^2$
admits a periodic tiling \citep[Theorem~1.1]{Bhattacharya}.  Greenfeld and Tao subsequently
gave a quantitative version of the planar result and proved that the
translation set of every tiling of $\Z^2$ by one finite tile is a disjoint
union of finitely many
sets, each invariant under some nonzero translation \citep[Theorems~1.3\textup{(i)} and~1.5]{GTStructure}.
An earlier planar result of \citet[Theorem~5.4]{WijshoffVanLeeuwen} concerns
polyominoes without holes and produces a tiling whose translation set
is a single lattice. The unrestricted finite-tile theorem allows
arbitrary disconnected sets and only asks for a finite-index period
subgroup. For Euclidean tilings, \citet[Theorem~1.1]{RationalPolygons}
prove periodic-tiling existence for rational polygonal sets, including
disconnected sets and sets with holes. These positive planar results
do not establish the conjecture for arbitrary bounded measurable
Euclidean tiles. We use the lattice results only for the least-dimension
conclusion above; the Euclidean assertion of \cref{thm:main} is proved
directly below.

\citet{GTPeriodic} disproved the periodic tiling conjecture in sufficiently
large dimension.  Their construction first produces a tile in
$\Z^2\times G_0$, with $G_0$ a finite abelian $2$-group, by encoding a
Sudoku rule with an arithmetic obstruction to periodicity.  The present
proof follows this general strategy and the $p$-adic Sudoku formulation of
\citet[Section~4]{GTUndecidable}.  We also use the idea of combining
simultaneous tiling equations into one equation from
\citet[Section~3]{GTPeriodic}.  The quotient-to-lattice reduction of
\citet[Theorem~1.1 and Corollary~1.2]{MSSReduction} explains the relevance of controlling the finite
factor: $\Z^2$ times a finite cyclic group is a quotient of $\Z^3$.

Here the Sudoku object is an array $W(n,m)$ over a finite alphabet,
with $n$ in a finite interval $I$ of columns and $m\in\Z$ indexing rows.
A finite set of allowed words prescribes the possible samples
$(W(n,dn+e))_{n\in I}$ along every integer-slope line. The arithmetic
word rule has two features: one such array exists with every column
nonconstant, but no array with those properties has a vertical period.
We turn these finite word constraints and column conditions into
simultaneous tiling equations.

There are two requirements at each stage of this strategy. The encoding
must admit at least one tiling, and every fully periodic tiling of the
encoded object must yield a fully periodic solution of the preceding
constraints. A system of tiling equations therefore means several
equations with the \emph{same} unknown translation set. Separate
solutions of the equations would not suffice. The finite group matters
as well: a general finite abelian group needs several generators, whereas
a cyclic group can be represented by just the third integer coordinate.

The central construction here encodes the Sudoku constraints while
keeping that finite factor cyclic. For each word position, a finite-valued
function has one input digit designated for each symbol. A symbol is
called active when changing that digit can change the function's value.
An arbitrary solution may activate several symbols at one position;
the constraints are imposed on every resulting choice of symbols.
The explicit solution we construct activates just one.
Cycle tests exclude forbidden simultaneous dependencies, and activation
tests force at least one symbol at every ordinary position.  Two seed
channels force distinct values in every column of the resulting Sudoku
array.  The activation tests use a nonuniform list of output shifts with
uniform coordinate counts on a product of two prime-order groups.
These counts force activity, while suitable permutations of the
remaining coordinates give an explicit common solution to all the
equations. Both seed tests use the same auxiliary output function.  All finite group factors have pairwise coprime
orders, so their product remains cyclic. The horizontal residue pairs
needed by the seed tests only index values of one cyclic output
coordinate; they do not become an additional group factor.

A stacking construction using a fresh prime then gives one tile in
$\Z^2\times\Zmod{Q}$.  We supply a direct rigidification argument that
passes to $\Z^3$ and also handles the unit-cube thickening in $\R^3$.
Integer vertices alone would not justify restricting an arbitrary
Euclidean tiling to integer translations. Our construction uses two
component shapes with the same large common part and one displaced
marked cube. The common part forces the component centers in a
hypothetical periodic tiling onto a translated grid. Coverage of the
marked cubes then forces invariance in the quotient kernel, allowing
that tiling to descend.  Thus the Euclidean argument includes arbitrary real
translation vectors.  The result concerns full periodicity; neither
connectedness of the tile nor absence of every individual nonzero
period is asserted.

\Cref{fig:proof-route} records the two directions maintained by these
constructions. The proof is organized as follows.  \Cref{sec:word} proves the Sudoku
obstruction and gives its last-nonzero-digit model.
\Cref{sec:encoding} constructs the cyclic group and the dependence tests;
\cref{sec:activation} forces activity and rules out periodic common
solutions.  \Cref{sec:model} constructs one common solution explicitly.
\Cref{sec:stacking,sec:geometry} perform stacking and the passage to three
dimensions.  Every construction and nonperiodicity argument needed for
\cref{thm:main} is proved in the paper.

\begin{figure}[!htbp]
\centering
\begin{tikzpicture}[
  every node/.style={font=\small,align=center},
  stage/.style={draw=MidnightBlue!65,fill=MidnightBlue!4,
    rounded corners=2pt,text width=3.15cm,minimum height=1.75cm,inner sep=4pt},
  route/.style={-{Latex[length=2mm]},line width=.7pt}
]
\node[stage] (word) {Word arrays\\with nonconstant\\columns\\\cref{sec:word}};
\node[stage,right=5mm of word] (system) {Common solution\\in $\Z^2\times V$\\$V$ cyclic\\\cref{sec:encoding,sec:activation,sec:model}};
\node[stage,right=5mm of system] (stack) {One tile $F$\\in $\Z^2\times\Zmod Q$\\\cref{sec:stacking}};
\node[stage,right=5mm of stack] (space) {One finite tile $T$\\and its thickening\\$T+[0,1]^3$\\\cref{sec:geometry}};
\draw[route] (word) -- (system);
\draw[route] (system) -- (stack);
\draw[route] (stack) -- (space);
\draw[route,MidnightBlue] ([yshift=4mm]word.north west) --
  node[above=1pt] {construct a solution, then a tiling}
  ([yshift=4mm]space.north east);
\draw[route,BurntOrange] ([yshift=-4mm]space.south east) --
  node[below=1pt] {a fully periodic tiling would give a forbidden vertical period}
  ([yshift=-4mm]word.south west);
\end{tikzpicture}
\caption{The two directions of the proof. The construction proceeds to
 the right, starting from the explicit last-nonzero-digit array in
 \cref{word:model}. To exclude full periodicity, any hypothetical periodic tiling
 is decoded in the reverse direction, producing a line-rule array
 with nonconstant columns and a vertical period. Every tiling equation at the
 second stage uses the same translation set, and the last stage uses
 the same $T$ in both the discrete and Euclidean assertions.}
\label{fig:proof-route}
\end{figure}

